\documentclass[11pt]{article}

% Essential packages (matched to the cipher-maps companion preamble)
\usepackage[margin=1in]{geometry}
\usepackage{amsmath,amssymb,amsthm}
\usepackage{mathtools}
\usepackage[numbers,square]{natbib}
\usepackage{booktabs}
\usepackage{hyperref}
\hypersetup{hypertexnames=false}
\usepackage{cleveref}
\usepackage{algorithm}
\usepackage{algpseudocode}

% Theorem environments
\theoremstyle{definition}
\newtheorem{definition}{Definition}[section]

\theoremstyle{plain}
\newtheorem{theorem}{Theorem}[section]
\newtheorem{lemma}[theorem]{Lemma}
\newtheorem{corollary}[theorem]{Corollary}
\newtheorem{proposition}[theorem]{Proposition}

\theoremstyle{remark}
\newtheorem{remark}{Remark}[section]
\newtheorem{example}{Example}[section]

% Notation (aligned with the cipher-maps companion where shared)
\newcommand{\enc}{\mathrm{enc}}
\newcommand{\dec}{\mathrm{dec}}
\newcommand{\TV}{d_{\mathrm{TV}}}
\newcommand{\TVsamp}{d_{\mathrm{TV},\,\mathrm{samp}}} % sampling-noise TV floor
\newcommand{\B}{\{0,1\}}
\newcommand{\GF}{\mathrm{GF}(2)}
\newcommand{\spn}{\mathrm{span}}
\newcommand{\class}{\mathrm{class}}
\DeclareMathOperator{\Uniform}{Uniform}
\DeclareMathOperator{\rank}{rank}

% Additional notation for this paper
\newcommand{\W}{W}                          % stored-pattern span
\newcommand{\Rz}{\mathcal{R}(z)}            % row space of the solution matrix
\newcommand{\proj}{\pi}                     % class-naming quotient projection
\newcommand{\Quot}{Q}                       % class quotient GF(2)^M / C
\newcommand{\supp}{\mathrm{supp}}           % stored value support
\newcommand{\Adv}{\mathsf{Adv}}             % distinguishing advantage
\newcommand{\al}{\alpha}                    % codespace share alpha(v)
\newcommand{\bz}{\mathbf{z}}                % solution matrix
\newcommand{\ay}{a_y}                       % band-indicator vector for key y
\newcommand{\cx}{c(x)}                      % stored canonical codeword of x
\newcommand{\lawreal}{\mu_{\mathrm{real}}}  % real non-member output law
\newcommand{\lawM}{\mu_{M1}}                % idealized (M1) non-member law
\newcommand{\inter}{\cap}

\title{Codec-Controlled Retrieval:\\
Structural Frequency-Hiding from the Non-Member Channel of XOR Retrieval}
\author{Alexander Towell\\\texttt{lex@metafunctor.com}}
\date{June 2026}

\begin{document}

\maketitle

\begin{abstract}
A retrieval data structure (Bloomier, ribbon, XOR, binary-fuse) returns a value
for \emph{every} key, with no membership gate. For keys it was built on it returns
the stored value in near-optimal space; for every other key it returns, in the
words of the literature, an arbitrary value, junk, a don't-care. We show that
junk is a designable object. For homogeneous ribbon (XOR) retrieval, where a query
is the $\GF$-linear form $\mathrm{lookup}(y)=\ay^{\mathsf T}\bz$, the value returned
for a non-member key lies in the $\GF$ span $\W$ of the stored codewords, and under
an idealized uniform-on-$\W$ model it follows a law fixed by a public value codec
rather than by the private data. That law is provably independent of how often each
value was stored: the span depends only on \emph{which} values are present, never on
their multiplicities. The control is governed by a sharp threshold. Full codec
control holds if and only if the stored codewords are transversal to the codec's
class partition, equivalently $\rank \proj|_{\W} = \log_2 K$, a step function of an
integer rank with no graded regime; below it whole codec classes receive exactly
zero probability. We cash this out as a structural frequency-hiding property: in a
distinguishing game (FreqDist) an adversary with non-member oracle access, even with
unboundedly many adaptive queries, has advantage exactly zero in the idealized model
(the two oracles are the identical law) and, in the real construction, the two
non-member laws are within total variation $\delta(p_0)+\delta(p_1)$ of each other, so
the advantage is at most that distance, where $\delta$ is a small, scale-independent,
redundancy-governed deviation we measure to be at or below the query-sampling floor
($0.00334$ at $10^5$ queries) up to $10^7$ keys. We are honest about the perimeter:
a second, multi-instance attack (a coincidence oracle with accuracy
$1-\tfrac12\sum_v \al(v)^t$, defeated by \emph{concentrated} rather than flat codespace
shares, so concentration is the better defense there) pulls the codec the opposite way,
so there is no single optimal codec, only a public tunable frontier. The construction is a header-only
C++23 implementation, and every theorem has a passing empirical check at scale. The
result is the concrete $\GF$-linear instantiation of a frequency-hiding property
that online encrypted-database defenses pay for per query.
\end{abstract}

% =====================================================================
% Source of record: source/construction-note.md (the proven, reviewed
% maph note). This draft RESHAPES that content into PoPETs venue prose;
% all theorem statements, bounds, and numbers are transcribed faithfully.
% =====================================================================

\section{Introduction}
\label{sec:intro}

Outsourced and encrypted key-value lookup has a persistent weakness: it leaks how
often each value occurs. Deterministic and property-preserving encryption, used to
make an encrypted store searchable, is the cautionary baseline. Equal ciphertexts
mean equal plaintexts, so the histogram of stored ciphertexts is the histogram of
plaintexts up to relabeling, and an adversary with auxiliary frequency knowledge
recovers the mapping. This is the inference attack of Naveed, Kamara, and Wright on
property-preserving encrypted databases~\cite{naveed2015inference}. The defenses
that answer it, frequency-smoothing encryption~\cite{lacharite2018fse} and
PANCAKE~\cite{grubbs2020pancake}, hide frequency by acting \emph{online}: they pad,
re-encode, or inject fake accesses at query or insert time, paying bandwidth or
maintenance per operation.

This paper studies a surface those defenses overlook. A \emph{retrieval} structure,
the static-function family that includes Bloomier filters~\cite{chazelle2004bloomier},
ribbon and BuRR retrieval~\cite{dillinger2021ribbon,dillinger2022burr}, and
binary-fuse filters~\cite{graf2022binaryfuse,graf2020xor}, returns a value for every
key, member or not, with no membership test. For a key in the built set it returns
the stored value in space close to the information-theoretic
minimum~\cite{dietzfelbinger2008succinct}. For a key not
in the set it returns something the literature is unanimous in calling
\emph{arbitrary}: a don't-care, the price of a construction that only needs to be
correct on the member set. XOR and fuse filters do \emph{specify} how unused
\emph{slots} are filled (uniform-random or zero) as an engineering choice, but that is
distinct from characterizing the decoded \emph{value} distribution on non-member
\emph{keys}, which to our knowledge the retrieval literature does not study: as recently
as 2025, the state of the art in static retrieval~\cite{hu2025retrieval} leaves the
off-set output distribution uncharacterized.

We fill exactly that silence, and find it was hiding a designable, provable object.
Our vehicle is homogeneous ribbon retrieval, in which a query is a $\GF$-linear
functional of a fixed solution matrix: $\mathrm{lookup}(y) = \ay^{\mathsf T}\bz$,
where $\bz$ stacks the solved solution rows and $\ay$ is the per-key coefficient
vector. (We stress at the outset that $\ay$ is a \emph{dense} band of roughly
$64$ bits, not a sparse $3$-wise XOR; the linear-algebra view, not the
sparse-hypergraph view, is what drives our results.) Because a query is linear, its
reachable outputs live in a subspace, and that single fact organizes everything that
follows.

\paragraph{Contributions.}
We state the contribution in three honestly-separated tiers, leading with what is
new and crediting what is not.
\begin{itemize}
  \item \textbf{(A, baseline, cited not claimed.)} That a value codec's codespace
  allocation can shape a non-member output distribution is \emph{known}. It is the
  author's own result for a random-oracle singular-hash-map
  construction~\cite{towell_bernoulli_maps}, and it is the
  distribution-transforming-encoder idea behind Honey
  Encryption~\cite{juels2014honey,cheng2019pmte}. We use it as established
  background and do not headline it.
  \item \textbf{(B, novel.)} The same codec-output law \emph{survives} the move to a
  $\GF$-linear retrieval structure, where the non-member output is an XOR of solution
  rows confined to a stored subspace, not a fresh uniform hash. Uniformity is no
  longer assumed; it becomes a statement about a row space, and it can fail. We prove
  the output support is exactly the stored span (\cref{thm:support}), give the
  idealized law (\cref{thm:idealized}), and prove it is independent of storage
  frequency (\cref{thm:freqindep}).
  \item \textbf{(C, headline.)} Control holds or fails by a \emph{sharp threshold}.
  The non-member law equals the codec's codespace shares if and only if the stored
  codewords are transversal to the codec's class partition, $\rank \proj|_{\W} =
  \log_2 K$, a step function of an integer rank with no graded regime
  (\cref{thm:threshold}). Below it, whole codec classes receive exactly zero
  probability. We are not aware of a prior statement of this threshold in the retrieval
  or filter literature; its underlying counting step (a subspace meets each coset it
  touches in a constant number of points, \cref{lem:subspace-coset}) is the coset weight
  distribution / coset-counting fact of coding theory~\cite{macwilliams1977}, which we
  apply in a new setting (the non-member output of a linear retrieval structure) to
  obtain a security-relevant codec-control rule.
\end{itemize}
On the security side, we formalize a distinguishing game (FreqDist) in which the
idealized adversary advantage is \emph{exactly zero} (\cref{thm:freqdist-ideal}) and, in
the real construction, the two non-member laws are within total variation
$\delta(p_0)+\delta(p_1)$ of each other, so the advantage of an adversary issuing any
number of adaptive non-member queries is at most that distance (\cref{thm:freqdist-real}),
where $\delta$ is the small, scale-independent deviation of the real construction from the
idealization, measured to be at or below the query-sampling floor up to ten million keys. We implement the construction
(header-only C++23) and validate every theorem empirically at scale. Finally, we are
honest about the perimeter: a structurally different multi-instance attack, a
coincidence oracle, is defeated by the \emph{opposite} codec choice, so there is no
single optimal codec, only a public tunable lever between two security objectives.

\paragraph{Threat model, in one breath.}
The adversary has oracle access to a single deployed instance and queries it on keys
of its choice, observing returned values. It restricts attention to \emph{non-member}
keys (member keys trivially return their stored value), and its goal is to recover
the storage-frequency distribution of values in the member set. We defend exactly
this channel against exactly this attack; access-pattern, volume, and size leakage
are out of scope, and the static structure offers no online guarantee. The honest
statement is narrow and strong.

\paragraph{Positioning.}
This is the construction-and-experiments companion to the cipher-maps
framework~\cite{towell2026ciphermaps}, whose abstraction has random-oracle and
RecSplit backends but no linear-retrieval backend. We supply the $\GF$-linear backend
and prove its structural frequency-hiding property. The opposite-pull between our two
attacks is, concretely, the linear-retrieval instance of the cipher-maps
$(\TV, \text{length})$ Pareto frontier.

\section{Background and the linear-retrieval view}
\label{sec:prelim}

We give the $\GF$ lens and the notation the theorems need. Throughout, arithmetic on
$M$-bit patterns is over the field $\GF$: vector addition is bitwise XOR, scalar
multiplication is bitwise AND, and an $M$-bit pattern is an element of the vector
space $\GF^M$. We write $a^{\mathsf T}\bz$ for the $\GF$ matrix-vector product (an
XOR of selected rows of $\bz$), $v \in \W$ for membership, $A \subseteq B$ for
inclusion, $A \inter B$ for intersection, and $|A|$ for cardinality.

\subsection{Codecs, classes, and codespace shares}

A value codec is a pair of maps $\enc : V \to \GF^M$ and $\dec : \GF^M \to V$ over an
$M$-bit pattern space, with $\dec(\enc(v)) = v$ for every logical value $v$. The
codec partitions $\GF^M$ into \emph{decode classes}
\[
  \class(v) = \{\, p \in \GF^M : \dec(p) = v \,\},
\]
one per value, and assigns each value a canonical pattern $\enc(v)$ in its class. The
\emph{codespace share} of $v$ is the fraction of patterns that decode to it,
\[
  \al(v) = \frac{|\class(v)|}{2^M}.
\]
A \emph{balanced} codec gives every one of its $K$ values a codeword of equal length
$\log_2 K$, so every class has the same size and $\al(v) = 1/K$. A \emph{skewed} codec
(a genuine Huffman code, say) gives shorter codewords to favored values, so classes
have unequal sizes and $\al$ is concentrated. We use \texttt{prefix\_codec}, whose
canonical patterns are \emph{left-aligned}: $\enc(v)$ places $v$'s codeword in the top
codeword bits and zeros the low bits. This left-alignment is a structural feature we
exploit later.

\subsection{Ribbon retrieval as a linear form}

A homogeneous ribbon retrieval structure over a key set $S$ stores a solution vector
of $m$ entries, each an $M$-bit pattern; stack them as the rows of an $m \times M$
matrix $\bz$ over $\GF$. For a key $y$ the structure computes a band specification
$(\mathrm{start}, \mathrm{coeffs})$, where $\mathrm{coeffs}$ is a $64$-bit odd word
and $\mathrm{start}$ is a legal band offset, and returns the XOR of the solution rows
selected by the set bits of $\mathrm{coeffs}$ placed at offset $\mathrm{start}$.
Collecting those selected positions into a band-indicator vector $\ay \in \B^m$ (bit
$\mathrm{start}+j$ of $\ay$ equals bit $j$ of $\mathrm{coeffs}$, and $0$ elsewhere),
\begin{equation}
  \mathrm{lookup}(y) = \ay^{\mathsf T}\bz.
  \label{eq:linear}
\end{equation}
This is the central structural fact: \emph{lookup is a $\GF$-linear function of the
fixed solution matrix $\bz$}, with the per-key coefficient vector $\ay$. The vector
$\ay$ is a dense band of up to $64$ set bits (about $32$ on average, since a
pseudorandom odd word has about half its bits set), all confined to a single
$64$-slot window. It is emphatically \emph{not} a $3$-wise XOR. We use only linearity
in this section.

The builder solves the banded $\GF$ system, by forward elimination plus
back-substitution over the short random band of each row~\cite{dietzfelbinger2019gauss},
so that, for every member $x \in S$,
\begin{equation}
  \ay[x]^{\mathsf T}\bz \;=\; \cx, \qquad \cx = \enc(v(x)),
  \label{eq:member}
\end{equation}
where $\cx$ is the stored canonical codeword of $x$'s value (the right-hand side for
$x$'s row); ``solvable instance'' means forward elimination plus back-substitution
produced a $\bz$ satisfying every such equation, which the builder verifies before
returning. Ribbon retrieval realizes this in space near $r|S|$ bits for an $r$-bit
value, riding on a practical, well-studied substrate~\cite{dillinger2022burr};
nothing in our analysis disturbs that space bound.

\subsection{The stored span and the idealized model}

Define the \emph{stored-pattern span}
\begin{equation}
  \W \;=\; \spn_{\GF}\{\, \cx : x \in S \,\},
  \label{eq:Wdef}
\end{equation}
the $\GF$ subspace of $\GF^M$ generated by the stored canonical codeword patterns. We
also write
\[
  \Rz = \{\, a^{\mathsf T}\bz : a \in \B^m \,\} = \spn_{\GF}\{\text{rows of }\bz\}
\]
for the row space of $\bz$, the set of all $\GF$ combinations of its rows. By
\eqref{eq:linear} every query output is an element of $\Rz$, so the reachable outputs
are confined to a subspace fixed by what was stored. The whole paper is the
consequence of that confinement.

Two of our results characterize the \emph{real} construction (the output support
$\W$, and the deviation $\delta$); two characterize an \emph{idealized} model, which
we name and flag up front.
\begin{definition}[Idealized model M1]
\label{def:M1}
For a non-member query $y$, model the pre-decoding output as a uniform draw from the
stored span:
\[
  \text{M1:}\qquad \ay^{\mathsf T}\bz \;\sim\; \Uniform(\W),
  \quad\text{i.e.}\quad
  \Pr[\ay^{\mathsf T}\bz = w] = 1/|\W| \text{ for every } w \in \W.
\]
\end{definition}
M1 is an \emph{idealization}, not a derived fact: it posits flatness across the
proven support $\W$. The real band $\ay$ is the deterministic image of a fingerprint
of $y$, its $64$-bit word is pinned odd, and a single band cannot reach two rows more
than $63$ apart, so the real band induces only a \emph{near}-uniform law on $\W$.
We adopt M1 as the maximally-mixing reference, study its consequence in
\cref{sec:t1t2,sec:t3,sec:t4}, and quantify the gap to the real construction
separately in \cref{sec:t5}. Keeping that division explicit is what lets us be precise
about which claims are theorems and which are measurements.

\section{The non-member output is uniform on the stored span}
\label{sec:t1t2}

This section establishes the support of the non-member output (a proven fact about
the real construction) and the idealized law on that support. Throughout we fix one
solvable ribbon instance over $S$ with a value codec, and read the builder's behavior
as stated in \cref{sec:prelim}. We separate the two results: \cref{thm:support}
(support) is unconditional; \cref{thm:idealized} (law) is conditional on M1.

\subsection{The output support is exactly the stored span}

\begin{theorem}[Output support]
\label{thm:support}
The row space of the solved solution matrix satisfies $\Rz = \W$. Consequently every
query output (member or non-member) lies in $\W$, and the reachable non-member outputs
span $\W$.
\end{theorem}

\begin{proof}[Proof sketch]
We prove the two inclusions; the full induction is in \cref{app:proofs}.
\emph{($\W \subseteq \Rz$.)} Each generator $\cx$ of $\W$ is a member output:
$\cx = \ay[x]^{\mathsf T}\bz$ by \eqref{eq:member}, hence a row-space element by
\eqref{eq:linear}. A subspace generated by elements of $\Rz$ lies in $\Rz$, so
$\W \subseteq \Rz$.
\emph{($\Rz \subseteq \W$.)} We show every row of $\bz$ lies in $\W$, by reading how
back-substitution fills $\bz$. The builder value-initializes every slot to the zero
pattern and assigns only pivot columns, so (P1) free slots are zero (and $0 \in \W$).
For pivot slots, the right-hand side carried through forward elimination is always an
XOR of the stored values $\cx$ and of earlier such right-hand sides, hence an element
of $\W$; back-substitution then sets each pivot slot to such an element XORed with
already-assigned downstream slots, which by induction on decreasing column index are
themselves in $\W$. Out-of-range references contribute $0 \in \W$. Hence (P2) every
pivot slot lies in $\W$. By (P1) and (P2) every row of $\bz$ is in $\W$, so a query
output $a^{\mathsf T}\bz$, an XOR of rows, is in $\W$; thus $\Rz \subseteq \W$.
Combining, $\Rz = \W$. Every output therefore lies in $\W$, and since each generator
$\cx$ is realized by its own member key $x$, the reachable non-member outputs span
$\W$.
\end{proof}

\Cref{thm:support} is the proven invariant: the row space, quantified over
\emph{all} coefficient vectors $a \in \B^m$ (not only the realizable band-indicators
$\ay$), equals $\W$. We claim only what is true about the reachable set: it is
contained in $\W$ and spans $\W$, but it need not equal $\W$ as a raw set, since a
non-generator element of $\W$ (an XOR of two or more $\cx$) need not be the output of
any single band-indicator.

\begin{remark}[Genericity, stated explicitly]
\label{rem:genericity}
\Cref{thm:support} uses exactly two facts about the solved instance: free slots are
zero, and every occupied slot is a $\GF$ combination of the stored right-hand sides.
Both are guaranteed by the builder as written: no step injects a pattern outside $\W$
(no random fill of free slots, no fingerprint salt added to a value). A future builder
variant that filled free slots with nonzero patterns, or added a non-$\W$ constant to
a right-hand side, would break this; the property is one of \emph{this} builder, not
of ribbon retrieval in the abstract. A computational check (tag \texttt{[span]})
confirms the inclusion: across $30{,}000$ distinct non-member queries every output's
re-encoded canonical pattern landed in $\W$, with zero violations.
\end{remark}

\subsection{The idealized law on the stored span}

\begin{theorem}[Idealized non-member law]
\label{thm:idealized}
Under the idealized model M1 (\cref{def:M1}), the non-member output law is
\begin{equation}
  \Pr[\mathrm{out}(y) = v] \;=\; \frac{|\class(v) \inter \W|}{|\W|}.
  \label{eq:idealized}
\end{equation}
\end{theorem}

\begin{proof}
The codec composition decodes the pre-decoding output, so
$\mathrm{out}(y) = \dec(\ay^{\mathsf T}\bz)$, and therefore
\[
  \Pr[\mathrm{out}(y) = v] = \Pr[\ay^{\mathsf T}\bz \in \class(v)].
\]
Under M1 the pre-decoding output is uniform on $\W$, so the probability of landing in
any subset of $\GF^M$ equals the fraction of $\W$ that the subset contains:
\[
  \Pr[\ay^{\mathsf T}\bz \in \class(v)]
  = \Pr[\ay^{\mathsf T}\bz \in \class(v) \inter \W]
  = \frac{|\class(v) \inter \W|}{|\W|}.
\]
The intersection with $\W$ is automatic, since $\ay^{\mathsf T}\bz \in \W$ always by
\cref{thm:support}. A value whose entire class misses $\W$ receives probability zero,
regardless of how the codec weights it; this geometric fact is the seed of the
threshold in \cref{sec:t4}.
\end{proof}

\Cref{thm:idealized} is exactly the conditional statement ``given the uniform-on-$\W$
action of M1, the codec-controlled formula holds''. It says nothing about how close
the real construction comes to M1; the real band is only near-uniform on $\W$, and
quantifying that residual is the separate subject of \cref{sec:t5}.

\begin{remark}[Why we adopt M1 on the image, not on the band support]
One might hope to derive M1 from the more primitive assumption that $\ay$ is uniform
on its band support $A$. The standard tool (image of a uniform vector under a linear
map: $L(U)$ is uniform on $L(A)$ when fibers are equal-sized) requires $A$ to be a
subspace surjecting onto $\W$. It does not apply here: $A$ excludes the zero vector,
pins the low bit of the band word to $1$, and confines all set bits to a single
window, so $A$ is not a subspace and the equal-fiber hypothesis is generically false.
Surjectivity $L(A) = \W$ alone does not give uniformity. We therefore adopt uniformity
on the image $\W$ directly. (For the linear map itself, $\mathrm{im}\, L = \Rz = \W$
over the \emph{full} domain $\B^m$, which is \cref{thm:support}; we do not conflate
that with the stronger and false claim that the band support alone surjects with equal
fibers.)
\end{remark}

\section{Frequency independence}
\label{sec:t3}

We now isolate exactly what the non-member law depends on. Storing a key-value map
fixes two separate things: the value \emph{support} (the set of distinct values that
occur) and the \emph{multiplicities} (how many keys carry each value). The next result
says the non-member law is a function of the support and the codec only; the
multiplicities never enter. This is the security lemma the frequency-analysis result
of \cref{sec:security} rests on.

Fix the codec. A key-value assignment over $S$ induces, for each value $v$ in the
support, a stored canonical pattern $\cx = \enc(v(x))$ that depends only on $v(x)$,
not on $x$: two keys mapped to the same value store the same pattern. Write the
multiplicities $m_v = |\{x \in S : v(x) = v\}|$ and the underlying \emph{set} of
distinct stored patterns
\[
  \supp = \{\, \cx : x \in S \,\} = \{\, \enc(v) : v \in \text{support} \,\}.
\]
The set $\supp$ has one element per distinct value, regardless of how many keys
realize it; it is fixed by the value support alone.

\begin{theorem}[Frequency independence]
\label{thm:freqindep}
Fix a codec. Let two key-value assignments have the same value support but arbitrary,
possibly very different, multiplicities $m_v$ and $m_v'$. Then they induce the same
stored span $\W$, and therefore the same idealized non-member law
$\Pr[\mathrm{out} = v] = |\class(v) \inter \W| / |\W|$.
\end{theorem}

\begin{proof}
By \eqref{eq:Wdef}, $\W = \spn_{\GF}\{\cx : x \in S\} = \spn_{\GF}(\supp)$. The span of
a multiset of vectors equals the span of its underlying set: a repeated generator adds
nothing to the set of $\GF$ combinations, since using it twice cancels
($g \oplus g = 0$) and using it once is already available from a single copy. Hence
$\W = \spn_{\GF}(\supp)$ depends on $\supp$ only, and $\supp$ is fixed by the value
support independently of the $m_v$. So $\W$ is invariant under any change of the $m_v$
that preserves the support. By \cref{thm:idealized} the idealized law is a function of
$\W$ and the codec's decode classes only, both independent of the $m_v$; therefore the
idealized non-member law is independent of the storage frequencies.
\end{proof}

In the real construction the two laws are not exactly equal: each build solves a
different banded system and so has its own near-uniform (rather than exactly uniform)
law on the shared $\W$. The two real laws therefore differ only by the per-build
deviation from $\Uniform(\W)$ characterized in \cref{sec:t5}, a quantity bounded there
and independent of the storage frequencies; no frequency-dependent term appears in
either the idealized law or that error budget.

\paragraph{Consequence for security.}
An adversary who probes the structure with non-member keys is sampling from the law of
\eqref{eq:idealized} (up to the \cref{sec:t5} deviation). Because the multiplicities
$m_v$ never enter that law, no amount of non-member probing reveals them: two databases
with identical support but opposite frequency profiles (one almost entirely a single
value, the other balanced) present statistically indistinguishable non-member outputs.
The channel leaks the support and the codec, both public design choices, and nothing
about how often each value is stored. This is what lets the FreqDist game of
\cref{sec:security} argue that classical frequency analysis has no purchase here.

\paragraph{Computational confirmation.}
The contrastive check (tag \texttt{[contrastive]}) builds two structures with the same
codec (prefix lengths $A{:}1$, $B{:}2$, $C{:}2$) and the same support $\{A,B,C\}$ but
storage frequencies $99\% / 0.5\% / 0.5\%$ and $50\% / 25\% / 25\%$, a raw frequency
gap up to $0.49$. Over $30{,}000$ distinct keys the two non-member distributions agree
per value to within $0.0035$ ($A$), $0.0044$ ($B$), $0.0009$ ($C$), far below the
$0.49$ gap that would appear if storage frequency leaked, and both sit within $0.03$ of
the codec shares $0.50 / 0.25 / 0.25$ predicted by \cref{thm:idealized}. The residual
is the per-build deviation of \cref{sec:t5}, not a frequency-dependent term.

\section{The sharp codec-control threshold}
\label{sec:t4}

\Cref{thm:idealized} gives the idealized law $|\class(v) \inter \W| / |\W|$;
\cref{thm:freqindep} shows it depends only on the support and the codec. This section
answers the next question precisely: for which supports does that law equal the
codec's own codespace shares, $\Pr[\mathrm{out} = v] = \al(v)$ for \emph{every} $v$
(``full codec control''), and what happens at the boundary? The answer is a single
sharp threshold, stated as transversality of $\W$ to the codec's class partition, with
\emph{no} intermediate regime. This is the headline result.

\subsection{A balanced linear codec and its class partition}

We work with a balanced linear codec: all $K = 2^{\log_2 K}$ values have codewords of
equal length $\log_2 K$, so $\al(v) = 1/K$ for every $v$. ``Linear'' is a property of
the class partition, made precise as follows. Let the \emph{homophone subspace} be the
set of patterns whose top $\log_2 K$ codeword bits are zero,
\[
  C = \{\, p \in \GF^M : \text{the top } \log_2 K \text{ bits of } p \text{ are } 0 \,\},
\]
a $\GF$ subspace of dimension $M - \log_2 K$ with $|C| = 2^{M - \log_2 K}$. The key
structural fact is that the codec's classes are exactly the cosets of $C$: two
patterns decode equally iff they share their top $\log_2 K$ bits, iff they differ by an
element of $C$. So each class is a coset $t + C$ of size $|C|$, and
$\al(v) = |C|/2^M = 2^{-\log_2 K} = 1/K$. Form the value quotient
\[
  \Quot = \GF^M / C,
\]
a $\GF$ space of dimension $\log_2 K$ whose $K$ points are the $K$ classes, and let
$\proj : \GF^M \to \Quot$ be the quotient map (projection onto the top $\log_2 K$
codeword bits). Then $\proj(p)$ names the class of $p$, and $\proj$ is $\GF$-linear and
surjective with kernel $C$. For \texttt{prefix\_codec} the canonical patterns are
left-aligned, so they lie in a transversal of $C$ and $\proj$ is injective on the set
of canonical patterns; we use this at the end.

\subsection{The subspace-meets-coset lemma}

\begin{lemma}[Subspace meets coset]
\label{lem:subspace-coset}
Let $\W$ be a $\GF$ subspace of $\GF^M$ and $t + C$ any coset of $C$. Then
$\W \inter (t + C)$ is either empty or a coset of $\W \inter C$; in particular its size
is either $0$ or $|\W \inter C|$, a constant independent of $t$.
\end{lemma}

\begin{proof}
If the intersection is empty there is nothing to prove. Otherwise pick
$w_0 \in \W \inter (t + C)$; we claim $\W \inter (t + C) = w_0 + (\W \inter C)$. For
the forward inclusion, take $w \in \W \inter (t + C)$. Both $w, w_0 \in \W$, so
$w \oplus w_0 \in \W$; both lie in $t + C$, so $w \oplus w_0 \in
(t+C) \oplus (t+C) = C$. Hence $w \oplus w_0 \in \W \inter C$ and
$w = w_0 \oplus (w \oplus w_0) \in w_0 + (\W \inter C)$. For the reverse, any
$w_0 \oplus u$ with $u \in \W \inter C$ lies in $\W$ (both terms in $\W$) and in
$t + C$ ($w_0 \in t+C$, $u \in C$). Thus the intersection is the coset
$w_0 + (\W \inter C)$, of size $|\W \inter C|$.
\end{proof}

The lemma is the crux: a subspace meets every class it touches in the \emph{same}
number of points, so the per-class mass is two-valued, a fixed positive constant on hit
classes and zero on missed classes. There is no room for a class to be partially filled
at an intermediate density.\footnote{This constant coset-hit count is folklore linear
algebra over $\GF$, equivalent to the coset weight distribution / coset-counting fact of
coding theory~\cite{macwilliams1977}. We claim no novelty for the lemma itself; the new
content is reading it as the non-member output law of a retrieval structure and the
control threshold (\cref{thm:threshold}) it yields.}

\subsection{\texorpdfstring{The mass on each class is $0$ or $1/K'$}{The mass on each class is 0 or 1/K'}}

Apply \cref{lem:subspace-coset} with the classes $\class(v) = t_v + C$. The per-class
idealized mass is $q(v) = |\class(v) \inter \W| / |\W|$, which is $0$ if $\W$ misses
$\class(v)$ and $|\W \inter C| / |\W|$ (the same value on every hit class) otherwise. A
class is hit iff some $w \in \W$ has $\proj(w) = \proj(t_v)$, that is, iff the class,
as a point of $\Quot$, lies in $\proj(\W)$. So the number of hit classes is
\[
  K' := |\proj(\W)| = 2^{\rank \proj|_{\W}},
\]
a subspace of $\Quot$, hence a power of two. Restricting $\proj$ to $\W$, the map
$\proj|_{\W} : \W \to \Quot$ has image $\proj(\W)$ and kernel $\W \inter C$, so the
first isomorphism theorem (equivalently rank-nullity over $\GF$) gives the counting
identity $|\W| = |\proj(\W)| \cdot |\W \inter C| = K' \cdot |\W \inter C|$. The constant
hit-class mass is therefore
\[
  \frac{|\W \inter C|}{|\W|} = \frac{1}{|\proj(\W)|} = \frac{1}{K'}.
\]
The entire idealized law is the step function
\[
  q(v) =
  \begin{cases}
    1/K' & \text{if class } v \text{ is hit } (\proj(t_v) \in \proj(\W)), \\
    0 & \text{if class } v \text{ is missed},
  \end{cases}
\]
with exactly $K'$ classes hit, summing to $K' \cdot (1/K') = 1$.

\subsection{Full control iff transversality, and sharpness}

\begin{theorem}[Sharp codec-control threshold]
\label{thm:threshold}
For a balanced linear codec, full codec control ($q(v) = \al(v) = 1/K$ for every $v$)
holds if and only if $\W$ is transversal to the class partition, that is
\[
  \proj(\W) = \Quot \quad\Longleftrightarrow\quad \rank \proj|_{\W} = \log_2 K.
\]
The transition is sharp: when $\proj(\W)$ is a proper subspace of $\Quot$ ($K' < K$,
i.e.\ $\rank \proj|_{\W} < \log_2 K$), the $K'$ hit classes are each over-weighted to
$1/K' > 1/K$ by the exact factor $K/K' > 1$, and the $K - K'$ missed classes are
exactly zero. There is no intermediate regime: the mass vector is a step function of
the integer $\rank \proj|_{\W}$, flipping to the flat $1/K$ profile precisely at
$\rank \proj|_{\W} = \log_2 K$.
\end{theorem}

\begin{proof}
Full control requires $q(v) = 1/K$ for every $v$, hence (from the step function) every
class hit \emph{and} the common hit-class mass equal to $1/K$, i.e.\ $1/K' = 1/K$, i.e.\
$K' = K$. Both hold iff $\proj(\W) = \Quot$, equivalently $\rank \proj|_{\W} = \log_2 K$
(this is the transversality condition: $\W$ meets every coset of $C$). For sharpness,
$K' = 2^{\rank \proj|_{\W}}$ is a power of two; raising $K'$ to the next attainable
value is the only way the mass vector can change, and it jumps discontinuously. Below
threshold the hit classes carry $1/K' > 1/K$ exactly and the missed classes carry $0$
exactly, with no value strictly between $0$ and $1/K$ and no gradual drift toward
$\al$.
\end{proof}

\subsection{The engineering bridge and a worked instance}

For \texttt{prefix\_codec} the threshold becomes a concrete, checkable condition. The
canonical patterns are left-aligned, so they lie in the transversal
$T = \{p : \text{low } M - \log_2 K \text{ bits zero}\}$, itself a subspace with
$T \inter C = \{0\}$ (an element of $T$ has its low bits zero, an element of $C$ has its
top bits zero, and only $0$ satisfies both). Since the stored patterns lie in $T$ and
$T$ is a subspace, $\W \subseteq T$, hence $\W \inter C \subseteq T \inter C = \{0\}$,
so $\W \inter C = \{0\}$. With trivial kernel, $\proj$ is injective on all of $\W$, so
$|\proj(\W)| = |\W|$ and
\begin{equation}
  \rank \proj|_{\W} = \dim \W = \mathrm{gf2\_rank}(\text{stored canonical patterns}).
  \label{eq:bridge}
\end{equation}
The transversality threshold $\rank \proj|_{\W} = \log_2 K$ therefore reads as the
engineering rule
\[
  \mathrm{gf2\_rank}(\text{stored canonical patterns}) = \log_2 K:
\]
store at least $\log_2 K$ linearly independent canonical patterns (enough distinct
values that the stored codewords span the full codeword space). Below that rank,
control is lost in the sharp way above; at or above it (the rank saturates at $\log_2 K$,
since $\proj(\W)$ is a subspace of the $\log_2 K$-dimensional $\Quot$), control is full.

\begin{example}[The length-2 codec on $M = 4$, $K = 4$]
\label{ex:m4}
The canonical patterns are $A = 0000$, $B = 0100$, $C = 1000$, $D = 1100$, each class of
size $4$, $\al = 1/4$, with $\log_2 K = 2$.
\begin{itemize}
  \item Store $\{A, B\}$: patterns $\{0000, 0100\}$, $\mathrm{gf2\_rank} = 1 < 2$. Then
  $\W = \{0000, 0100\}$, $K' = 2$, and the predicted mass over $[A,B,C,D]$ is
  $[\tfrac12, \tfrac12, 0, 0]$: the two hit classes over-weighted to $1/K' = 1/2$, the
  rest exactly zero.
  \item Store $\{A, B, C\}$: patterns $\{0000, 0100, 1000\}$, $\mathrm{gf2\_rank} = 2 =
  \log_2 K$. Now $D = 1100 = 0100 \oplus 1000$ lies in $\W$ \emph{though it was never
  stored}, so $\W = \{0000, 0100, 1000, 1100\}$, $\proj(\W) = \Quot$, $K' = 4$, and the
  predicted mass is $[\tfrac14, \tfrac14, \tfrac14, \tfrac14] = \al$. This is full
  control, and it confirms that what matters is membership of a class in $\proj(\W)$, not
  whether that class's value was stored.
\end{itemize}
The \texttt{[threshold]} check asserts these two vectors exactly (tolerance $10^{-9}$,
the masses being dyadic rationals over the enumerated $|\W|$), with $\mathrm{gf2\_rank}$
pinned at $1$ and $2$ respectively, tying the flip to the rank crossing $\log_2 K$. (The
hit-class mass is $1/K' = |\W \inter C|/|\W|$ in general; it coincides with $1/|\W|$ here
only because left-alignment forces $\W \inter C = \{0\}$, so $|\W \inter C| = 1$.)
\end{example}

\begin{remark}[The general skewed codec is open]
\label{rem:skewed}
The coset-counting argument uses one feature of the balanced case essentially: a
\emph{single} homophone subspace $C$ describes the whole partition, so
\cref{lem:subspace-coset} applies uniformly. A skewed codec (variable codeword lengths)
has classes of unequal sizes $2^{M - \ell_v}$ that are cosets of \emph{different}
subspaces $C_v$, so there is no common $C$ and the two-valued conclusion fails: different
classes can in principle carry different positive densities. Full control still has a
clean abstract statement, $|\class(v) \inter \W|/|\W| = \al(v)$ for every $v$ ($\W$ in
general position with respect to the unequal partition), but characterizing which
supports achieve it, and whether the skewed transition is still sharp or genuinely
graded, is \emph{not} settled by the balanced argument, and we do not claim it. We state
the balanced threshold as the proven, sharp result and explicitly defer the skewed case
as open.
\end{remark}

\section{From the idealized model to the real construction}
\label{sec:t5}

\Cref{thm:idealized} is an idealization (M1): the real band is only near-uniform on
$\W$, for three reasons recorded in \cref{def:M1}. We close that gap here, honestly.
Writing $\lawreal$ for the actual non-member output law of a solved build and $\lawM$
for the $\Uniform(\W)$-induced law, the object of interest is
\[
  \delta \;:=\; \TV(\lawreal, \lawM).
\]
Above the span threshold (\cref{thm:threshold}) the M1 law coincides with the codec's
codespace shares $\al$, so $\delta$ is the total-variation distance between the real
non-member law and $\al$. We treat this as a \emph{characterization} result: we attempt
a clean closed-form bound, find it resists, and deliver an empirical characterization
together with an honestly-labeled big-$O$ statement. We do \emph{not} prove a
closed-form bound on $\delta$.

\paragraph{Small and scale-independent.}
The scale experiment (E3, \cref{sec:experiments}) fixes one above-threshold
configuration, the balanced $M = 8$, $K = 8$ codec storing all eight values (stored
patterns span a rank-$3 = \log_2 8$ subspace at every $N$), and sweeps the stored-key
count $N$ across four orders of magnitude, measuring $\delta$ as mean
TV-to-codespace over build seeds with $n_q = 10^5$ non-member queries per build. The
deviation is \emph{small} (order $0.003$ in TV, two orders of magnitude below the
$0.1$ scale that would signal a breakdown) and \emph{scale-independent} (flat in $N$
from $10^4$ to $10^7$, if anything drifting weakly downward, with all confidence
intervals overlapping). The ribbon's $\epsilon$ auto-scaling, which widens the band as
$N$ grows to keep large systems solvable, does not degrade codec control: near-uniformity
on $\W$ holds at ten million keys as at ten thousand.

\paragraph{The deviation is at or below the sampling floor.}
The crux of an honest characterization is to decide whether the measured $\sim 0.003$ is
the \emph{systematic} deviation of the true law from M1 or merely the query
sampling-noise floor. Each build estimates the non-member law from $n_q$ finite queries
over $K$ classes; even if the true law were exactly uniform-on-$\W$ ($\delta = 0$), the
empirical frequencies would fluctuate, producing a strictly positive observed
TV-to-truth purely from sampling. Quantify that floor. Under M1 above threshold each
class has probability $1/K$, the observed share is a $\mathrm{Binomial}(n_q, p_v)/n_q$
count, and the expected sampling TV is
\[
  \mathbb{E}[\TVsamp] = \tfrac12 \sum_v \mathbb{E}|\hat p_v - p_v|.
\]
The summand is a \emph{mean absolute deviation} (not a standard error); for an
approximately normal $\hat p_v$ it is the factor $\sqrt{2/\pi} \approx 0.80$ smaller than
the standard error $\sqrt{p_v(1 - p_v)/n_q}$. Evaluating for the E3 settings ($K = 8$
equal classes, $p_v = 1/8$, $n_q = 10^5$) via de Moivre's exact binomial MAD identity
(the normal form agrees to four decimals),
\[
  \mathbb{E}[\TVsamp] = 0.00334.
\]
(A reader who instead sums per-class standard errors gets the cruder upper proxy
$0.00418$, exactly the $1/0.80 = 1.25$ factor larger; we use the correct MAD-based
$0.00334$.) The measured E3 mean TV is $0.0028$ to $0.0035$, and the floor $0.00334$ lies
\emph{inside} all four confidence intervals, with the $N = 10^6$ point essentially on it
and the $N = 10^7$ mean ($0.00284$) below it. The measured deviation therefore
\emph{straddles} the floor rather than exceeding it. Modeling the systematic part $D$ as
adding in quadrature with the sampling noise $S$ (heuristic, since TV is an
$L^1$/mean-absolute quantity, not $L^2$), the most over-floor point implies
$D \le 0.0011$ and the at-or-below points imply $D$ indistinguishable from zero. The
strongest honest statement: at $n_q = 10^5$ the systematic deviation of the real band
from $\Uniform(\W)$ is at or below the $\sim 0.0033$ sampling floor, below measurement
resolution for support comfortably above threshold.

\paragraph{The one visible effect is redundancy, not scale.}
The genuine, visible deviation is governed not by $N$ but by storage
\emph{redundancy}: how far above the span threshold $\W$ sits, i.e.\ how many independent
stored keys back each codec class. The frequency-independence experiment (E2,
\cref{sec:experiments}) fixes the skewed length-$\{1,2,2\}$ codec on support
$\{V_0, V_1, V_2\}$, all profiles at rank $2 = \log_2 K$, and reports a per-build
same-profile cross-seed baseline (build one profile twice, TV the two non-member laws),
isolating per-build variance plus sampling noise. The three rich profiles (uniform
$33/33/33$, mid $50/25/25$, skew\_b $90/5/5$) sit at the floor ($\sim 0.003$); the thin
skew\_a profile ($99/0.5/0.5$) is elevated to $\sim 0.019$. The cause is thin support:
at $99/0.5/0.5$ the minority values are each stored only about $25$ times, so although
$\W$ reaches rank $2$, it does so with almost no redundancy, sitting at the threshold
edge, and the solved band on that barely-spanning $\W$ is meaningfully less uniform per
build. So $\delta$ is a function of the per-class storage redundancy and is essentially
independent of $N$ for fixed redundancy. The scale-independence of E3 and the
redundancy-dependence of E2 are the same statement from two directions: growing $N$ at
fixed support richness does not move $\delta$, but thinning a class toward the threshold
does.

\begin{remark}[The elevation does not undermine frequency independence]
\label{rem:elevation}
Even in the thin-support elevated regime, the elevated deviation does \emph{not} encode
storage frequency. The cross-profile advantage of E2 tracks the per-build baseline of
the noisier profile, not the frequency gap: the three pairs that include skew\_a all land
at the same $\sim 0.014$ magnitude despite very different frequency gaps, of the same
order as skew\_a's own $\sim 0.019$ baseline, while a real $33$-vs-$90$ gap between two
\emph{rich} profiles sits at the $\sim 0.0036$ floor. The elevation is symmetric
per-build noise of the thin profile, not a frequency-difference signal: it loosens the
constant in the FreqDist bound (\cref{sec:security}) without introducing any term that
grows with the frequency gap.
\end{remark}

\begin{remark}[Honestly labeled deliverable]
\label{rem:t5-bound}
We state the empirical characterization, not a theorem. \textbf{Observation.} For
support comfortably above threshold (each class backed by many independent keys),
$\delta = O(\text{sampling floor})$ and is indistinguishable from zero at the measured
query budget, scale-independently in $N$ up to $10^7$. As $\W$ approaches the threshold
(a class backed by a handful of keys), $\delta$ grows mildly, to $\sim 0.017$ at the
edge (minority classes stored $\sim 25$ times), governed by the per-class store count,
not by $N$ or by the storage frequencies. \textbf{Conjecture (qualitative, not proven).}
The visible part of $\delta$ is controlled by the redundancy beyond the threshold
dimension, equivalently by the minimum per-class store count $m_{\min}$; $\delta$ sits at
the floor once $m_{\min}$ is large and degrades only as $m_{\min}$ falls toward the
minimum needed to span. A structural proof, presumably bounding the non-uniformity of a
banded $\GF$ solution on $\W$ in terms of the number of independent equations pinning
$\W$, is left open. This is deliberately weaker than \cref{thm:support,thm:idealized,%
thm:freqindep,thm:threshold}, which are proven; \cref{sec:t5} is where the real
construction meets the idealization, and the honest deliverable is the measurement plus
the big-$O$.
\end{remark}

\section{Security: frequency-analysis resistance}
\label{sec:security}

We now cash out the theory as a security property. \Cref{thm:freqindep} established the
lemma it rests on (the non-member law depends on the support and the codec only, never
the multiplicities); \cref{sec:t5} established the deviation budget the real construction
spends. This section states the game, proves the idealized advantage is exactly zero,
bounds the real advantage by the \cref{sec:t5} budget, and is careful about the
perimeter.

\subsection{Threat model and the FreqDist game}

The adversary has \emph{oracle} access to a single deployed instance: it queries chosen
keys and observes returned values, adaptively. The restriction that makes for a clean
statement is that it queries only \emph{non-member} keys (member keys return their stored
canonical value by construction; the interesting channel is what the structure says about
keys it was not built on). Its goal is to learn the storage-frequency distribution among
the member set, the multiplicities $m_v$. This is the classical frequency-analysis
surface: a deterministic substitution cipher, or deterministic encryption used as an
encrypted-search index, leaks plaintext frequency directly, the attack of Naveed, Kamara,
and Wright~\cite{naveed2015inference} that PANCAKE~\cite{grubbs2020pancake} neutralizes
by smoothing the access distribution. Our point is that the non-member channel here leaks
frequency \emph{structurally not at all}: not because frequencies are smoothed or padded,
but because they never enter the law the channel samples from (\cref{thm:freqindep}).

\begin{definition}[The FreqDist game]
\label{def:freqdist}
Fix a value support $\supp$ and a codec, both \emph{public}.
\begin{enumerate}
  \item The adversary $\mathcal{A}$ chooses two storage-frequency distributions $p_0$
  and $p_1$ over the \emph{same} support $\supp$, with arbitrary, possibly very
  different, multiplicities (e.g.\ $p_0$ piles almost all mass on one value, $p_1$ is
  balanced; only the set of distinct values must agree, which fixes a common $\W$ by
  \cref{thm:freqindep}).
  \item The challenger samples $b \leftarrow \{0,1\}$, builds the structure from a key
  set realizing $p_b$, and gives $\mathcal{A}$ oracle access.
  \item $\mathcal{A}$ issues non-member queries, observes the decoded outputs, and outputs
  a guess $b'$.
\end{enumerate}
The advantage is $\Adv(\mathcal{A}) = |2\Pr[b' = b] - 1|$, the standard distinguishing
advantage. The game is meaningful above the span threshold (\cref{thm:threshold}); we
assume $\rank \proj|_{\W} = \log_2 K$, the natural deployment regime in which both
profiles induce full codec control over the common $\W$.
\end{definition}

\subsection{Idealized advantage is exactly zero}

\begin{theorem}[FreqDist, idealized]
\label{thm:freqdist-ideal}
Under the idealized model M1 the FreqDist advantage is exactly zero:
$\Adv(\mathcal{A}) = 0$ for every adversary, computationally unbounded and adaptive
included.
\end{theorem}

\begin{proof}
By \cref{thm:freqindep}, $\W$ is a function of $\supp$ alone and invariant under any
change of the multiplicities preserving the support; $p_0$ and $p_1$ share $\supp$, hence
induce the same $\W$. By \cref{thm:idealized} the idealized non-member law under either
profile is $|\class(v) \inter \W|/|\W|$, a function of $\W$ and the codec's classes only,
with no multiplicity term. Both inputs are identical for $p_0$ and $p_1$, so the
non-member law under $p_0$ equals that under $p_1$: the two oracles are the \emph{same}
distribution. No function of identically distributed observations can distinguish their
source, so $\Pr[b' = b] = 1/2$ for every $\mathcal{A}$ regardless of the number or
adaptivity of queries, and $\Adv(\mathcal{A}) = 0$.
\end{proof}

The advantage is zero not because distinguishing is hard but because there is nothing to
distinguish: the two oracles are one and the same law, and identical distributions cannot
be told apart at any sample size, by any number of adaptive queries.

\subsection{The two real laws are close, so the advantage is small at any query count}

In the real construction the two oracles are not exactly equal: each profile's build
solves its own banded system, inducing its own near-uniform law on the shared $\W$. Write
$\lawreal(p_b)$ for the real non-member law under $p_b$ and recall
$\delta(p_b) = \TV(\lawreal(p_b), \lawM)$, the per-build deviation from the shared M1 law
(the same object for both, being frequency-independent by \cref{thm:freqindep}). The
correct statement here is about the \emph{distance between the two real laws}, not a
per-query advantage cap: it must hold no matter how many non-member queries the adversary
draws, since \cref{def:freqdist} grants unboundedly many adaptive queries.

\begin{theorem}[FreqDist, real: the two laws are close]
\label{thm:freqdist-real}
The two real non-member laws are close in total variation, frequency-independently and
regardless of how many samples are drawn:
\[
  \TV\big(\lawreal(p_0), \lawreal(p_1)\big) \;\le\; \delta(p_0) + \delta(p_1).
\]
Consequently, in the convention $\Adv = |2\Pr[b'=b]-1|$, the per-query distinguishing
advantage (the advantage of the optimal test on one observed non-member output) is
\[
  \Adv(\mathcal{A}) \;\le\; \TV\big(\lawreal(p_0), \lawreal(p_1)\big)
  \;\le\; \delta(p_0) + \delta(p_1),
\]
and the only quantity an adversary can sharpen with more queries is its estimate of these
two \emph{fixed} laws, whose separation is capped by the same
$\delta(p_0)+\delta(p_1)$ at every query count and contains no frequency-dependent term.
\end{theorem}

\begin{proof}
Both real laws are within $\delta$ of the same idealized M1 law, which is
frequency-independent by \cref{thm:freqindep} ($\W$ depends only on the shared support).
Routing through that common law and using the triangle inequality for TV,
\[
  \TV(\lawreal(p_0), \lawreal(p_1))
  \le \TV(\lawreal(p_0), \lawM) + \TV(\lawM, \lawreal(p_1))
  = \delta(p_0) + \delta(p_1).
\]
This compares the two \emph{fixed} per-build laws and is therefore independent of how many
samples the adversary observes: it is a statement about the laws, not about a sample. For
the advantage, the optimal single-observation test distinguishing two distributions has
advantage exactly their total variation in the $\Adv = |2\Pr[b'=b]-1|$ convention (Le Cam's
two-point bound), so on one observed output the advantage is at most
$\TV(\lawreal(p_0), \lawreal(p_1)) \le \delta(p_0) + \delta(p_1)$. (No factor of $2$
appears: the earlier single-observation reading multiplied the per-sample TV by the
convention factor; the per-observation advantage in this convention is the TV itself.)
A many-query adversary draws i.i.d.\ outputs from one fixed instance and can estimate that
instance's law to arbitrary precision, but the only thing this lets it do is decide which of
the two \emph{known, fixed} laws it faces, and those two laws are themselves within
$\delta(p_0)+\delta(p_1)$ of each other by the display above. Crucially that separation does
not grow with the number of queries and (by \cref{thm:freqindep}) carries no
frequency-dependent component: the adversary can sharpen its picture of each law without ever
uncovering a frequency-gap signal, because none is present. The whole of what more queries
can buy is bounded by the gap between two laws that differ only by the redundancy-governed,
frequency-independent $\delta$.
\end{proof}

The bound caps the gap between the two real laws, not a per-query quantity that could grow
with the number of queries. Plugging in the \cref{sec:t5} numbers: with both profiles rich,
$\delta(p_0) \approx \delta(p_1) \approx 0.003$, so the advantage is at most
$0.003 + 0.003 = 0.006$; with one profile at the threshold edge (the skew\_a regime),
$\delta \approx 0.017$ on the thin side and $\approx 0.003$ on the rich side, so the
advantage is at most $0.017 + 0.003 = 0.02$. E2 (\cref{sec:experiments}) measures exactly
this quantity, the total variation between the two realized non-member laws, and confirms
not just the bound but its \emph{shape}: the measured cross-profile distance sits at the
per-build floor and does not grow with the frequency gap between $p_0$ and $p_1$, exactly
as $\delta(p_0)+\delta(p_1)$ predicts and the idealized $\Adv = 0$ demands.

\subsection{An independence result, not a comparative one}

The claim is an \emph{independence} statement: the non-member channel is independent of
storage frequency. The idealized form is the cleanest possible such statement (the two
oracles are the identical distribution, $\Adv = 0$); the real form weakens it only to a
small constant set by per-build deviation, with no frequency-gap term. It is \emph{not} a
comparative claim of the form ``codec $X$ leaks $c$ times less frequency than codec $Y$'',
and deliberately not phrased that way. A comparative, directional reading is exactly what
admitted the sign error in the multi-instance coincidence-oracle analysis
(\cref{sec:reconciliation}), where ``minimizes the coincidence sum'' was misread as
``defends better'' when the defender direction is the reverse. An indistinguishability
statement carries no direction to invert: $\Adv = 0$ means the channel reveals nothing
about frequency, full stop, and the real bound caps the residual symmetrically in
$\delta(p_0)$ and $\delta(p_1)$ with no sign attached to which profile is more
concentrated. What the channel \emph{does} carry is only the public codec shares
$(\al(v))$ and the support, both public design choices; never the private $m_v$.

\subsection{Scope and boundary}

This section defends exactly one channel against exactly one attack.
\begin{itemize}
  \item \textbf{Covered.} The value-frequency channel of non-member outputs: an adversary
  with non-member oracle access, with any number of adaptive queries, cannot learn the
  storage-frequency distribution of values in $S$. Idealized advantage zero (M1, from
  \cref{thm:idealized,thm:freqindep}); in the real construction the two non-member laws are
  within total variation $\delta(p_0)+\delta(p_1)$, so the advantage is at most that
  distance, a small constant governed by storage redundancy and independent of the
  frequency gap (\cref{sec:t5}, E2).
  \item \textbf{Not covered.} Nothing about access-pattern leakage (which keys are
  queried, and correlations across queries), volume or size leakage, or the multi-instance
  coincidence oracle. That oracle is a \emph{different} game (correlating $t$ redeployments
  to predict membership from output coincidence, not frequency from one instance); it is
  defeated by \emph{concentration} of $\al$, the opposite codec choice, and is treated in
  \cref{sec:reconciliation}.
  \item \textbf{Standing assumption.} Above-threshold support
  (\cref{thm:threshold}) throughout, so both profiles realize full codec control and the
  common-$\W$ argument applies.
\end{itemize}
We do not claim a general confidentiality result. The honest statement is narrow and
strong: on its one channel, against its one attack, in its deployment regime, frequency
analysis has no purchase, structurally, by \cref{thm:freqindep}, with a quantified
real-world residual from \cref{sec:t5}.

\section{The two attacks pull opposite ways}
\label{sec:reconciliation}

The security story so far has one attack in it, single-instance frequency analysis,
defeated by frequency independence. There is a second, structurally different attack on
the \emph{same} construction, the multi-instance coincidence oracle, and the codec choice
that defeats one is the worst against the other. We state both, correct a sign error that
has circulated in companion writeups, and frame the resulting tension as a clean finding:
there is no single most-confidential codec.

\paragraph{Attack 1 (single instance): frequency analysis, defeated by flat codecs.}
The adversary sees one deployed instance and wants the $m_v$, as in classical frequency
analysis against a substitution cipher. \Cref{thm:freqindep} is the answer: the
non-member law is a function of the support and the codec only, so any above-threshold
codec defeats single-instance frequency analysis outright (the contrastive test confirms
output gaps $\sim 0.004$ between $99/0.5/0.5$ and $50/25/25$ storage). What residual the
channel \emph{does} leak is the public codec shares $\al$ plus the support. If one worries
that even this leaks value \emph{importance} (a Huffman codec's short codeword announces
``this value is frequent''), the way to minimize what $\al$ reveals is to \emph{flatten}
it: a uniform $\al = 1/K$ carries no per-value ordering. So within attack 1 the pressure
on the codec is toward \emph{flat} shares.

\paragraph{Attack 2 (multi-instance): the coincidence oracle, defeated by concentrated
codecs.} The adversary sees $t$ instances that share the same $(S, f)$ and differ only in
per-instance non-member randomness (a shared-function redeployment). The coincidence
oracle predicts ``$k \in S$'' iff all $t$ observed outputs are equal. For a member, all
$t$ canonical encodings equal $f(k)$, so coincidence always holds; for a non-member, the
$t$ outputs are i.i.d.\ draws from $\al$ and coincide with probability $\sum_v \al(v)^t$.
At a balanced prior the attacker accuracy is
\begin{equation}
  \mathrm{acc}(t) = 1 - \tfrac12 \sum_v \al(v)^t.
  \label{eq:acc}
\end{equation}
Write $S(t) = \sum_v \al(v)^t$. The defender wants $\mathrm{acc}(t)$ \emph{small}, hence
$S(t)$ \emph{large}. For fixed $\sum_v \al(v) = 1$ and $t \ge 2$, $S(t)$ is maximized by a
\emph{concentrated} (Huffman-like) distribution ($S(t) \to 1$ as mass piles onto a few
values) and minimized by the \emph{uniform} distribution ($S(t) = K^{1-t}$, small).
Therefore concentrated $\al$ gives $\mathrm{acc}(t)$ near $1/2$, the best defense, and
uniform $\al$ gives $\mathrm{acc}(t)$ near $1$, the worst. Within attack 2 the pressure is
toward \emph{concentrated} shares, the exact opposite of attack 1.

\begin{remark}[Correcting a sign error]
\label{rem:sign}
An earlier reading claimed the uniform codec is the \emph{better} multi-instance defense,
on the grounds that it minimizes $S(t)$ and minimizing the coincidence sum looks
protective. That reverses the defender direction: $\mathrm{acc}(t) = 1 - S(t)/2$, so a
\emph{smaller} $S(t)$ means \emph{higher} attacker accuracy, a worse defense. The ``uniform
defends best'' claim is the direction-reversed reading and is wrong; concentrated codecs
defend best. The corrected E4 result (\cref{sec:experiments}) confirms this direction
independently by Monte Carlo.
\end{remark}

\paragraph{The synthesis.}
The two attacks pull the codec in opposite directions: attack 1 is best served by flat
$\al$ (it already loses to independence, and flat shares minimize importance leakage);
attack 2 is best served by concentrated $\al$ (it minimizes \eqref{eq:acc} by maximizing
$S(t)$, and a Huffman codec is separately space-optimal). \textbf{There is no single
most-confidential codec.} Which corner you want depends on the threat model: at $t = 1$
(no redeployment) the single-instance pressure dominates and flat is preferable; at large
$t$ (heavy redeployment) the coincidence oracle dominates and Huffman is preferable. This
is exactly the linear-retrieval instance of the cipher-maps $(\TV, \text{length})$ Pareto
frontier~\cite{towell2026ciphermaps}: leakage (TV of the output law from a
frequency-revealing target, minimized at the flat corner) versus cost (expected codeword
length, minimized at the Huffman corner), now with a security reading on the cost axis,
the Huffman corner being coincidence-optimal as well as space-optimal. The opposite pull is
not a flaw but the expected shape of a Pareto problem: codec-controlled retrieval exposes a
single, public, tunable lever ($\al$) trading single-instance importance leakage against
multi-instance coincidence resistance (and space).

\section{Evaluation}
\label{sec:experiments}

We make the theory concrete and reproducible. Each experiment is generated by an artifact
in the implementation (a header-only C++23 benchmark binary or a Python analysis script);
the numbers below are read directly from the committed result files. Nothing here
introduces a new claim: E1 is the empirical face of the threshold (\cref{thm:threshold}),
E2 of frequency independence (\cref{thm:freqindep}) and the redundancy effect
(\cref{sec:t5}), E3 of the scale-independence and sampling-floor analysis (\cref{sec:t5}),
and E4 of the reconciliation direction (\cref{sec:reconciliation}). The construction is
implemented in the \texttt{maph} research playground; build and analysis scripts are
provided for reproduction.

\subsection{E1: the span threshold is sharp}

The benchmark fixes a balanced codec, sweeps the number of distinct stored values (hence
the $\GF$ rank of the stored canonical patterns), and reports the mean TV of the
non-member law to the codec shares $\al$, averaged over build seeds. \Cref{thm:threshold}
predicts a step: above the threshold $\rank \proj|_{\W} = \log_2 K$ the law equals $\al$
(TV near $0$); one rank below, the law cannot reach $\al$ and TV is large.
\Cref{tab:e1} shows the two flanking ranks.

\begin{table}[t]
  \centering
  \caption{E1: the non-member-to-codespace TV across the span threshold. The transition
  is a cliff, not a ramp. Sub-threshold confidence intervals are \emph{degenerate}
  (deterministic in rank), so the transition is not merely sharp but deterministic; only
  at and above threshold does the residual acquire width (the \cref{sec:t5} near-uniformity
  gap over finite queries).}
  \label{tab:e1}
  \begin{tabular}{lcc}
    \toprule
    config & rank just below $\log_2 K$ & rank $= \log_2 K$ \\
           & (rank, mean TV) & (rank, mean TV) \\
    \midrule
    \texttt{balanced\_M4\_K4} & $1$, \quad $0.500$ & $2$, \quad $0.00277$ \\
    \texttt{balanced\_M8\_K8} & $2$, \quad $0.500$ & $3$, \quad $0.00505$ \\
    \bottomrule
  \end{tabular}
\end{table}

TV drops from $0.5$ to $\sim 0.003$ (the \cref{sec:t5} floor) as the rank crosses
$\log_2 K$, with no intermediate value, exactly the step function of \cref{thm:threshold}.
Below threshold the per-config TV is deterministic in the rank (the step-function masses
are fixed dyadic rationals on the enumerated $\W$), so the sub-threshold confidence
intervals are degenerate: for instance \texttt{balanced\_M8\_K8} reads $0.875, 0.750,
0.500$ at ranks $0, 1, 2$ with zero-width intervals, while at threshold (rank $3$) the
interval acquires width $[0.00455, 0.00556]$. A skewed config in the same data shows the
analogous drop to $0.00238$ at rank $2$, with a non-degenerate sub-threshold value
($0.251$ at rank $1$) because the unequal class sizes make the below-threshold mass a
non-trivial average rather than a single dyadic constant.

\subsection{E2: FreqDist indistinguishability and the redundancy effect}

The benchmark fixes the skewed length-$\{1,2,2\}$ codec (codespace shares
$0.5/0.25/0.25$) on support $\{V_0, V_1, V_2\}$, all profiles above threshold (rank
$2 = \log_2 K$), and reports a same-profile cross-seed \emph{baseline} (build one profile
twice, TV the two non-member laws, isolating per-build variance plus sampling noise) and
the cross-profile \emph{advantage} (TV between two different storage profiles).
\Cref{tab:e2} gives both. \Cref{thm:freqindep} predicts the cross-profile advantage does
not track the frequency gap; \cref{sec:t5} predicts the one profile at the threshold edge
(skew\_a) carries an elevated per-build deviation.

\begin{table}[t]
  \centering
  \caption{E2: same-profile baseline and cross-profile advantage. The advantage tracks the
  noisier (thinner) profile, not the frequency gap: every pair including the thin skew\_a
  lands at $\sim 0.014$ regardless of the gap, while a real $33$-vs-$90$ gap between two
  rich profiles sits at the $\sim 0.0036$ floor.}
  \label{tab:e2}
  \begin{tabular}{lcc}
    \toprule
    profile / pair & same-profile baseline & cross-profile advantage \\
    \midrule
    uniform $(33/33/33)$ & $0.00292$ & --- \\
    mid $(50/25/25)$     & $0.00360$ & --- \\
    skew\_b $(90/5/5)$   & $0.00381$ & --- \\
    skew\_a $(99/0.5/0.5)$ & $0.01897$ & --- \\
    \midrule
    uniform vs mid       & --- & $0.00324$ \\
    uniform vs skew\_b   & --- & $0.00364$ \\
    mid vs skew\_b       & --- & $0.00402$ \\
    uniform vs skew\_a   & --- & $0.01336$ \\
    mid vs skew\_a       & --- & $0.01433$ \\
    skew\_b vs skew\_a   & --- & $0.01458$ \\
    \bottomrule
  \end{tabular}
\end{table}

The three rich profiles sit at the $\sim 0.003$ floor in both columns. The cross-profile
advantage tracks the noisier profile of the pair: the three pairs including skew\_a all
land at $\sim 0.014$, of the same order as skew\_a's own $\sim 0.019$ baseline, despite
storage-frequency gaps differing enormously; the genuinely large gap between two rich
profiles (uniform-vs-skew\_b) sits at the floor. So the elevation is a per-build property
of thin support (skew\_a stores its minority values only $\sim 25$ times each, leaving $\W$
at the threshold edge), a \cref{sec:t5} effect, not a frequency-leakage signal. The
cross-profile column is exactly the total variation $\TV(\lawreal(p_0), \lawreal(p_1))$
between the two realized non-member laws, the quantity \cref{thm:freqdist-real} bounds by
$\delta(p_0)+\delta(p_1)$ and that caps the advantage of an adversary at any query count.
This jointly confirms the independence (the cross-profile distance does not grow with the
frequency gap) and the redundancy characterization (it is set by the per-class store count
of the thinner build), realizing the real bound
$\Adv \le \TV(\lawreal(p_0), \lawreal(p_1)) \le \delta(p_0)+\delta(p_1)$ of
\cref{thm:freqdist-real}.

\subsection{E3: the deviation is small and scale-independent}

The benchmark fixes one above-threshold configuration (balanced $M = 8$, $K = 8$, all
eight values stored, rank $3 = \log_2 8$ at every $N$) and sweeps $N$ across four orders of
magnitude, reporting mean TV-to-codespace over build seeds with $10^5$ non-member queries
per build. \Cref{tab:e3} gives the result.

\begin{table}[t]
  \centering
  \caption{E3: mean TV-to-codespace versus stored-key count $N$. The deviation is flat in
  $N$ from $10^4$ to $10^7$ (overlapping intervals), and straddles the analytic
  query-sampling floor $0.00334$ (inside every interval, below it at $N = 10^7$). The
  ribbon's $\epsilon$ auto-scaling at $10^7$ keys does not degrade control.}
  \label{tab:e3}
  \begin{tabular}{lccc}
    \toprule
    $N$ & reps & mean TV & $95\%$ CI \\
    \midrule
    $10^4$ & $10$ & $0.00352$ & $[0.00313, 0.00392]$ \\
    $10^5$ & $10$ & $0.00311$ & $[0.00262, 0.00359]$ \\
    $10^6$ & $5$  & $0.00352$ & $[0.00272, 0.00432]$ \\
    $10^7$ & $3$  & $0.00284$ & $[0.00230, 0.00339]$ \\
    \bottomrule
  \end{tabular}
\end{table}

The deviation is order $0.003$ in TV (two orders of magnitude below a $0.1$ breakdown
scale) and flat in $N$ with overlapping intervals. The analytic sampling floor for these
settings ($K = 8$, $p_v = 1/8$, $n_q = 10^5$) is $0.00334$ (exact binomial MAD, matching
the normal approximation), and the measured values straddle it: it lies inside all four
intervals, the $N = 10^6$ mean sits on it, and the $N = 10^7$ mean ($0.00284$) falls
below. The systematic deviation is therefore at or below measurement resolution at this
query budget, consistent with \cref{sec:t5}.

\subsection{E4: the coincidence-oracle direction}

The script computes, for $K = 8$ values and $t = 1, \dots, 8$ redeployments, the attacker
accuracy $\mathrm{acc}(t) = 1 - \tfrac12 \sum_v \al(v)^t$ of \eqref{eq:acc} in closed form
and confirms it by Monte Carlo ($10^5$ member $+ 10^5$ non-member trials per rep, $10$
reps), across codecs of increasing concentration. Higher accuracy means worse defense; the
corrected direction is that concentration \emph{helps}. \Cref{tab:e4} gives the slice at
$t = 4$.

\begin{table}[t]
  \centering
  \caption{E4: coincidence-oracle attacker accuracy at $t = 4$, by codec concentration.
  Accuracy falls monotonically as the codec concentrates: uniform is the unique worst
  defense, the most-concentrated codec the best. Monte Carlo matches the closed form in
  every cell (within a $95\%$ interval), validating $\mathrm{acc}(t) = 1 - S(t)/2$.}
  \label{tab:e4}
  \begin{tabular}{lccc}
    \toprule
    codec (max share) & $\mathrm{acc}$ closed & $\mathrm{acc}$ MC & MC matches? \\
    \midrule
    uniform $(0.125)$              & $0.99902$ & $0.99903$ & yes \\
    intermediate\_zipf $(\sim 0.34)$ & $0.98240$ & $0.98255$ & yes \\
    huffman $(0.50)$              & $0.96667$ & $0.96671$ & yes \\
    intermediate\_padded $(0.90)$ & $0.67195$ & $0.67185$ & yes \\
    \bottomrule
  \end{tabular}
\end{table}

Attacker accuracy falls monotonically as the codec concentrates: uniform (every share
$0.125$) is the unique worst defense at $\sim 0.999$; the dyadic Huffman codec (top share
$0.50$) improves it to $\sim 0.967$; the most-concentrated codec (one value holding $0.90$
of the codespace) is the best defense at $\sim 0.672$. The Monte Carlo matches the closed
form in every cell. This is the direction \cref{sec:reconciliation} states: a smaller
coincidence sum $S(t)$ means higher attacker accuracy, so the uniform codec (which
minimizes $S(t)$) is the worst, not the best, multi-instance defense. E4 is thus the
quantitative backing for the opposite-pull finding: the codec that is harmless against
single-instance frequency analysis (uniform, where the advantage is zero regardless) is the
worst against the multi-instance coincidence oracle, and the concentrated codec reverses
both rankings.

\section{Related work and novelty}
\label{sec:related}

We position the contribution honestly, including against the author's own prior work, so
the genuinely new content is not confused with results that already exist. We organize the
positioning as a three-tier ledger.

\paragraph{Tier A: not novel here (reused baseline).}
The abstract law that a codec's codespace allocation controls the non-member output
distribution, $\Pr[\mathrm{out} = v] \propto |\class(v)|/2^M$, is \emph{not} a contribution
of this paper. It already appears in two places, cited as established baseline. First, the
author's singular-hash-map work~\cite{towell_bernoulli_maps} derives exactly this formula
for a \emph{random-oracle} construction, where the non-member output is a fresh hash of the
query and its distribution over codec classes is immediate. Second, the idea of designing an
encoder so its output follows a target distribution regardless of message frequencies is the
distribution-transforming-encoder idea of Honey Encryption~\cite{juels2014honey}, refined by
probability-model-transforming encoders~\cite{cheng2019pmte}. We credit both and do not claim
the abstract law as new.

\paragraph{Tier B: novel (the law survives the move to a linear retrieval structure).}
The novel content of \cref{thm:support,thm:idealized,thm:freqindep} is that the \emph{same}
codec-output law survives the move from a random-oracle construction to a $\GF$-linear
retrieval structure. There the non-member output is not a fresh hash but $a^{\mathsf T}\bz$,
an XOR of selected rows of the solved matrix, a deterministic linear function of the fixed
stored data. A priori there is no reason such a structure's off-target output should follow
the codec shares at all: the output is constrained to the stored span $\W$, fixed by what was
stored, and the per-key coefficient is a structured band, not a uniform draw.
\Cref{thm:support} proves the support is exactly $\W$; \cref{thm:idealized} (under M1) and
\cref{thm:threshold} (the threshold) characterize when the law on $\W$ equals the codec
shares; \cref{thm:freqindep} proves the law depends only on the support. The retrieval and
static-function literature that ribbon retrieval descends from, Bloomier
filters~\cite{chazelle2004bloomier}, ribbon and BuRR~\cite{dillinger2021ribbon,dillinger2022burr},
and binary-fuse filters~\cite{graf2022binaryfuse,graf2020xor}, treats the non-member output as
arbitrary, uncharacterized junk: it is the price of a construction that need only be correct on
the member set. These filters do \emph{specify} the fill of unused \emph{slots}
(uniform-random or zero) as an engineering choice, but that is a statement about empty slots,
not a characterization of the decoded \emph{value} distribution on non-member \emph{keys},
which to our knowledge is not studied; as of 2025 the state of the art in static
retrieval~\cite{hu2025retrieval} leaves that off-set output distribution uncharacterized.
Characterizing that distribution, and showing it is the codec's own codespace law gated by the
span threshold, is the new content.

\paragraph{Tier C: headline novelty (the sharp span threshold).}
The headline is the sharp threshold of \cref{thm:threshold}: the non-member law equals the
codec shares exactly when the stored patterns span the full codeword space
($\rank \proj|_{\W} = \log_2 K$), and the transition is a step in the integer rank, stated as
transversality of $\W$ to the codec's class partition. We are not aware of a prior statement
of this threshold for retrieval data structures: it is absent from the random-oracle
treatments (Tier A, where there is no span to be transversal to), and absent from the
retrieval literature (Tier B, which does not study the non-member law at all). The nearest
external \emph{technique} is coding theory's coset weight distribution / coset
counting~\cite{macwilliams1977}: the subspace-meets-coset lemma (\cref{lem:subspace-coset}),
which gives the constant coset-hit count, is folklore linear algebra over $\GF$, not itself
new. What is new is its application, reading that constant count as the non-member
\emph{value} law of a linear retrieval structure, and the resulting all-or-nothing
$\rank \proj|_{\W} = \log_2 K$ control threshold with its $\mathrm{gf2\_rank}$ engineering
bridge and security consequence; to our knowledge the coset-weight machinery has not been used
this way. The resulting $K' = |\proj(\W)|$ step function is the technical core, and the cliff
in E1 is its measured face.

\paragraph{Differentiation from adjacent work.}
Four adjacencies deserve an explicit boundary. \textbf{(1) The author's
singular-hash-map}~\cite{towell_bernoulli_maps} contains the identical formula but for a
random-oracle construction; the transfer to a linear structure, and the discovery that the
transfer is gated by a provable span threshold, is the new content. \textbf{(2) The author's
cipher-maps paper}~\cite{towell2026ciphermaps} shares the frequency-hiding goal, but its
backends are random-oracle and RecSplit with no linear-retrieval backend, and its concern is the
query-marginal distribution (which keys are asked), not the non-member output channel of a
linear structure; this paper supplies the linear backend it lacks and proves the structural
property for it. \textbf{(3) Honey-Encryption DTE}~\cite{juels2014honey,cheng2019pmte}
\emph{assumes} a uniform seed delivered by a cipher and builds on that assumption; we do not
assume uniformity, but \emph{prove} near-uniformity on the stored span from the linear structure
itself (M1 as idealization, \cref{sec:t5} as the measured gap), gated by the threshold.
\textbf{(4) Frequency-smoothing databases}: PANCAKE~\cite{grubbs2020pancake},
FSE~\cite{lacharite2018fse}, and the NKW inference attack~\cite{naveed2015inference} that
motivates them share the goal but operate at a different layer, online mechanisms that smooth the
access or encoding distribution at query or insert time (fake accesses, padding). Our property is
static and structural: the non-member law of a built linear structure is frequency-independent by
construction (\cref{thm:freqindep}), with no online smoothing, no fake traffic, no padding.

\paragraph{External filter-privacy and volume-hiding formalisms.}
Two recent external formalisms are the nearest cousins. Filic, Paterson, Unnikrishnan, and
Virdia~\cite{filic2022adversarial} give an adversarial-correctness-and-privacy model for
probabilistic data structures, the closest external filter-confidentiality formalism; it covers
membership PDS, not the retrieval \emph{value}-distribution, so the object we characterize (the
non-member output law of a retrieval structure) is outside its scope. Patel, Persiano, Yeo, and
Yung~\cite{patel2019volumehiding} give volume-hiding structured encryption, the closest
structural cousin in spirit: shape an observable (there, volume) to a public target via a
structural, hashed construction rather than online padding. Our work shapes a different observable
(the non-member value distribution) to a public target (the codec shares) via a different
structural mechanism (the $\GF$ span), and contributes the sharp threshold governing when the
shaping takes hold, which has no analogue in the volume-hiding setting.

\paragraph{Leakage-abuse attacks and leakage suppression.}
The broader backdrop is the searchable/structured-encryption leakage lineage that motivates
hiding structural channels in the first place. Leakage-abuse attacks established that
side-channel leakage from encrypted search is exploitable in practice: Islam, Kuzu, and
Kantarcioglu~\cite{islam2012access} recovered queries from access-pattern leakage, and Cash,
Grubbs, Perry, and Ristenpart~\cite{cash2015leakage} gave count and access-pattern
query-recovery and plaintext-recovery attacks across a range of leakage profiles, alongside the
NKW frequency inference attack~\cite{naveed2015inference} on property-preserving encryption. In
response, leakage-suppression frameworks reduce or eliminate specific leakage channels by
construction: Kamara, Moataz, and Ohrimenko~\cite{kamara2018structured} give general compilers
that suppress the query/search pattern, the framework inside which volume-hiding structured
encryption~\cite{patel2019volumehiding} sits. Our property is complementary to this line and
narrower in scope: it defends one specific structural channel, the value-frequency of non-member
outputs of a static linear retrieval structure, by construction and at zero per-query cost,
rather than suppressing query- or access-pattern leakage. It is a structural defense for a
channel that the leakage-abuse line did not target, not a substitute for the access-pattern and
search-pattern suppression those frameworks provide.

\section{Discussion, limitations, and conclusion}
\label{sec:conclusion}

We set out to fill a unanimous silence in the retrieval literature, that the value returned for
a non-member key is ``arbitrary'', and found the silence was hiding a designable object. For
homogeneous ribbon retrieval the non-member output lives in the $\GF$ span of the stored
codewords (\cref{thm:support}); under an idealized uniform-on-span model it follows a law fixed
by a public codec (\cref{thm:idealized}); that law is provably independent of how often each
value was stored (\cref{thm:freqindep}); and codec control of it holds or fails by a sharp
$\GF$-rank threshold with no graded regime (\cref{thm:threshold}). The real construction realizes
the idealization to within the query-sampling floor for rich support, scale-independently to ten
million keys (\cref{sec:t5}). Read as security, this is structural frequency-hiding: a
distinguishing adversary with non-member oracle access, issuing any number of adaptive
queries, has advantage exactly zero in the idealized model (\cref{thm:freqdist-ideal}), and
in the real construction the two non-member laws are within total variation
$\delta(p_0)+\delta(p_1)$, capping the advantage at that distance
(\cref{thm:freqdist-real}), a property of the static bytes with zero per-query cost, where
PANCAKE and FSE pay bandwidth or re-encoding online.

\paragraph{Limitations.}
Three boundaries are explicit. The threshold theorem (\cref{thm:threshold}) is proven for the
\emph{balanced} linear codec; the skewed (variable-length) case lacks a single homophone subspace,
its abstract full-control condition is stated but whether its transition stays sharp or becomes
genuinely graded is open (\cref{rem:skewed}). The deviation $\delta$ (\cref{sec:t5}) is
\emph{characterized}, not bounded in closed form: we deliver the measurement and a labeled
big-$O$, and conjecture but do not prove that $\delta$ is governed by the minimum per-class store
count (\cref{rem:t5-bound}). And the security treatment defends exactly one channel against
exactly one attack: the non-member value-frequency channel against single-instance frequency
analysis, plus the reconciliation of the opposite-pulling multi-instance coincidence oracle;
access-pattern, volume, and size leakage are not addressed, and the structure is static.

\paragraph{Future work.}
The skewed-codec threshold is the most pressing open problem: a characterization of which supports
put $\W$ in general position with respect to an unequal class partition, and whether the resulting
transition is sharp. A structural closed-form bound on $\delta$ in terms of storage redundancy
would upgrade \cref{sec:t5} from a measurement to a theorem. A white-box or randomized-encoding
angle (filling free slots with a designed pattern rather than zero, which \cref{rem:genericity}
flags as the place the support could be enlarged on purpose) may let the construction realize
output laws beyond the stored span. More broadly, this is the construction-and-experiments
companion to the cipher-maps program~\cite{towell2026ciphermaps}: it supplies the $\GF$-linear
retrieval backend that abstraction lacks, and the opposite-pull between the two attacks is the
linear-retrieval instance of its $(\TV, \text{length})$ frontier.

\paragraph{Takeaway.}
The values an XOR retrieval structure returns for non-member keys are not arbitrary junk: they are
uniform over the $\GF$ span of the stored codewords, so a public codec, not the private data,
determines that distribution, exactly when the stored values span the codec's class quotient,
which hands a static data structure the frequency-hiding property that online encrypted-database
defenses pay for per query.

\appendix

\section{Full proof of the output support}
\label{app:proofs}

We give in full the induction sketched in \cref{thm:support}, that every row of the solved
solution matrix $\bz$ lies in the stored span $\W$, so $\Rz \subseteq \W$. Combined with
$\W \subseteq \Rz$ (each generator $\cx = \ay[x]^{\mathsf T}\bz$ is a row-space element), this
gives $\Rz = \W$.

Read the back-substitution block of the builder, which value-initializes the solution to the zero
pattern, skips free columns, and assigns each pivot column the carried right-hand side XORed with
already-solved downstream slots:
\begin{verbatim}
  solution.assign(num_rows, 0);              // (i)   all slots 0
  for (col = num_rows; col-- > 0; ) {
      if (pivot_coeffs[col] == 0) continue;  // (ii)  skip free slots
      val = pivot_value[col];                // (iii) RHS for this pivot
      ... val ^= solution[ref] ...           // (iv)  XOR solved downstream
      solution[col] = val;                   // (v)   assign pivot slot
  }
\end{verbatim}

\paragraph{(P1) Free slots are zero.}
A column is a pivot column iff $\texttt{pivot\_coeffs[col]} \ne 0$ after forward elimination. Line
(i) value-initializes every entry to the zero pattern, and only pivot columns are ever assigned
(line (v) runs only when (ii) does not skip). Hence every non-pivot (free) column holds the zero
pattern, and $0 \in \W$ trivially (the empty XOR of generators).

\paragraph{(P2) Every pivot slot lies in $\W$.}
During forward elimination the carried right-hand side $\texttt{pivot\_value[col]}$ is built only
from the stored values by the update $v \mathrel{\oplus}= \texttt{pivot\_value[col]}$, starting from
$v = \cx$ for the incoming row of some $x \in S$. By induction each $\texttt{pivot\_value[col]}$ is
an XOR of a subset of the $\{\cx\}$ and of earlier $\texttt{pivot\_value}$ entries (themselves XORs
of the $\{\cx\}$), hence an element of $\W$. Back-substitution then sets $\texttt{solution[col]}$ to
$\texttt{pivot\_value[col]}$ XORed with already-assigned downstream entries
$\texttt{solution[ref]}$. We induct on \emph{decreasing} column index: the loop runs from the
largest column down, and each referenced $\texttt{ref} = \texttt{col} + \mathrm{offset} + \mathrm{bit}
> \texttt{col}$ was assigned in an earlier iteration, hence (by the inductive hypothesis) already
lies in $\W$. Out-of-range references ($\texttt{ref} \ge \texttt{num\_rows}$) are dropped and
contribute $0 \in \W$. So each $\texttt{solution[col]}$ is an XOR of elements of $\W$, hence in $\W$.

By (P1) and (P2) every row of $\bz$ is in $\W$. A query output $a^{\mathsf T}\bz$ is an XOR of rows
of $\bz$, hence an XOR of elements of $\W$, hence in $\W$, for every coefficient vector
$a \in \B^m$. Therefore $\Rz \subseteq \W$, and with $\W \subseteq \Rz$ we conclude $\Rz = \W$. The
full quantification over all $a \in \B^m$ (not only the realizable band-indicators $\ay$) is what
makes the equality hold; the reachable set of single-band outputs is contained in $\W$ and spans it,
but need not equal it as a raw set. \qed

\bibliographystyle{plainnat}
\bibliography{refs}

\end{document}
